\documentclass[10pt,a4paper]{article}

\usepackage[margin=20mm]{geometry}
\usepackage{amsmath,amssymb,amsthm}

\theoremstyle{definition}
\newtheorem{definition}{Definition}
\newtheorem*{problem*}{Problem}

\setlength{\parindent}{0pt}
\setlength{\parskip}{5pt}

\begin{document}

\begin{center}
  {\Large\bfseries Linear Matroid Parity in NC}
\end{center}
\section{Definition Linear Matroid Parity }
We write $[m]:=\{1,2,\ldots,m\}$.

\begin{definition}[Matroid]
A \emph{matroid} is a pair $M=(E,\mathcal I)$, where $E$ is a finite
\emph{ground set} and $\mathcal I\subseteq 2^E$ is a family of subsets,
called the \emph{independent sets}, satisfying the following axioms:
\begin{enumerate}
  \item $\varnothing\in\mathcal I$;
  \item if $I\in\mathcal I$ and $J\subseteq I$, then $J\in\mathcal I$;
  \item if $I,J\in\mathcal I$ and $|I|<|J|$, then there exists
  $e\in J\setminus I$ such that $I\cup\{e\}\in\mathcal I$.
\end{enumerate}
A maximal independent subset of $E$ is called a \emph{basis} of $M$.
\end{definition}

\begin{definition}[Rank function]
Let $M=(E,\mathcal I)$ be a matroid. For every $X\subseteq E$, the
\emph{rank} of $X$ in $M$ is
\begin{equation}
  \operatorname{rk}_M(X)
  :=\max\bigl\{|I|:I\subseteq X,\ I\in\mathcal I\bigr\}.
\end{equation}
Thus, $X$ is independent if and only if
$\operatorname{rk}_M(X)=|X|$.
\end{definition}

\begin{definition}[Linear matroid]
Let $\mathbb F$ be a field and let
$A\in\mathbb F^{r\times n}$ be a matrix whose columns are indexed by a
finite set $E$ of cardinality $n$. The \emph{linear matroid represented by
$A$}, denoted $M(A)$, is the matroid with ground set $E$ in which
$X\subseteq E$ is independent if and only if the columns of $A$ indexed by
$X$ are linearly independent over $\mathbb F$.

For every $X\subseteq E$, its matroid rank is the ordinary matrix rank:
\begin{equation}
  \operatorname{rk}_{M(A)}(X)=\operatorname{rank}_{\mathbb F} A[:,X].
\end{equation}
A matroid is \emph{linear} or \emph{representable over $\mathbb F$} if it
is isomorphic to $M(A)$ for some matrix $A$ over $\mathbb F$.
\end{definition}

Let $\mathbb F$ be a fixed finite field of odd characteristic. We consider a
matrix
\begin{equation}
  A=(a_1,b_1,a_2,b_2,\ldots,a_m,b_m)\in \mathbb F^{r\times 2m},
\end{equation}
where the symbols $a_i,b_i$ denote both the corresponding columns and their
column indices. The ground set of $M(A)$ is therefore
\begin{equation}
  E=\{a_1,b_1,a_2,b_2,\ldots,a_m,b_m\},
\end{equation}
and it is partitioned into the prescribed pairs
\begin{equation}
  \ell_i:=\{a_i,b_i\},\qquad i\in[m].
\end{equation}
For $S\subseteq[m]$, write
\begin{equation}
  \ell(S):=\bigcup_{i\in S}\ell_i
  \qquad\text{and}\qquad
  A[S]:=A[:,\ell(S)].
\end{equation}

\begin{definition}[Matroid parity]
Let $M=(E,\mathcal I)$ be a matroid whose ground set is partitioned into
pairs
\begin{equation}
  E=\ell_1\mathbin{\dot\cup}\ell_2\mathbin{\dot\cup}\cdots
  \mathbin{\dot\cup}\ell_m,
  \qquad |\ell_i|=2.
\end{equation}
A set $S\subseteq[m]$ is a \emph{parity set} if the union
$\ell(S)=\bigcup_{i\in S}\ell_i$ is independent in $M$. The
\emph{matroid parity problem} asks for a parity set of maximum cardinality.
When $M$ is represented by a matrix over a field, this is the
\emph{linear matroid parity problem}.
\end{definition}

\begin{definition}[Linear matroid parity value]
For the represented instance above, a set $S\subseteq[m]$ is a parity set
if and only if
\begin{equation}
  \operatorname{rank}_{\mathbb F} A[S]=2|S|.
\end{equation}
The \emph{linear matroid parity value} of $A$ is
\begin{equation}
  \nu(A):=\max\bigl\{|S|:S\subseteq[m],\
  \operatorname{rank}_{\mathbb F}A[S]=2|S|\bigr\}.
\end{equation}
\end{definition}

\section{Definition $\mathsf{NC}$}
\begin{definition}[$\mathsf{NC}$]
A language belongs to $\mathsf{NC}$ if it is decided by a logspace-uniform
family of Boolean circuits of polynomial size and depth $(\log n)^{O(1)}$.
The class $\mathsf{RNC}$ is the corresponding bounded-error randomized class.
\end{definition}

\section{Current best result}

For every fixed finite field $\mathbb F$ of characteristic different from $2$, linear matroid
parity has a randomized parallel algorithm of polynomial size and depth
$O((\log N)^2)$, where $N$ is the input length. Equivalently,
\begin{equation}
  \mathsf{LMP}_{\mathbb F}\in\mathsf{RNC}^2.
\end{equation}
The open problem is to remove the randomness while retaining polynomial size
and polylogarithmic depth.

\section{Problem formalization}
For the fixed field $\mathbb F$, define the threshold language
\begin{equation}
  \mathsf{LMP}_{\mathbb F}
  :=\bigl\{(A,k):A\in\mathbb F^{r\times 2m},\ 0\leq k\leq m,\
  \nu(A)\geq k\bigr\}.
\end{equation}


\begin{problem*}[Linear matroid parity in $\mathsf{NC}$]
\textbf{Decision version.}
Prove that
\begin{equation}
  \mathsf{LMP}_{\mathbb F}\in\mathsf{NC}.
\end{equation}
Equivalently, construct a deterministic logspace-uniform family of
polynomial-size, polylogarithmic-depth circuits which, given $(A,k)$, decides
whether $A$ contains $k$ prescribed pairs whose $2k$ columns are linearly
independent.
\end{problem*}
\end{document}
